\subsection{Vòng đời file và commit đầu tiên}
\label{subsec:c22-lifecycle}

Git theo dõi mỗi file trong kho chứa thông qua bốn trạng thái của vòng đời file, như minh họa ở Hình~\ref{fig:c22-lifecycle}.
Hiểu đúng vòng đời này giúp sinh viên không bao giờ mất bài do nhầm lẫn giữa file chưa lưu vào vùng stage và file đã commit \cite{progit}.
Ba vùng dữ liệu tương ứng là thư mục làm việc, vùng stage (còn gọi là index) và kho chứa gồm các commit \cite{atlassiangit}.

\begin{description}
\item[Untracked] File mới mà Git chưa quản lý, ví dụ trang \verb|About.aspx| vừa tạo trong project.
File untracked không xuất hiện trong lịch sử commit cho đến khi được đưa vào stage lần đầu.
\item[Modified] File đã tracked nhưng bị sửa đổi trong thư mục làm việc mà chưa đưa vào stage.
\item[Staged] File đã được \verb|git add| và nằm trong index, sẵn sàng cho commit kế tiếp.
\item[Committed] Dữ liệu đã nằm an toàn trong kho chứa dưới dạng một commit có mã băm SHA-1 duy nhất.
\end{description}

\begin{figure}[H]
\centering
\begin{tikzpicture}[>=Stealth, thick, box/.style={draw, rounded corners, fill=blue!10, minimum width=2.2cm, minimum height=0.9cm, align=center, font=\small}]
\node[box] (a) {Untracked};
\node[box, right=0.8cm of a] (b) {Modified};
\node[box, right=0.8cm of b] (c) {Staged};
\node[box, right=0.8cm of c] (d) {Committed};
\draw[->] (a) -- node[midway, above, font=\footnotesize] {add} (b);
\draw[->] (b) -- node[midway, above, font=\footnotesize] {add} (c);
\draw[->] (c) -- node[midway, above, font=\footnotesize] {commit} (d);
\draw[->, bend right=30] (d) to node[midway, below, font=\footnotesize] {edit} (b);
\end{tikzpicture}
\caption{Vòng đời của file trong Git từ chưa theo dõi đến đã commit}
\label{fig:c22-lifecycle}
\end{figure}

Lệnh \verb|git status| cho biết trạng thái hiện tại của từng file, còn \verb|git add| chuyển file sang vùng stage.
Dạng rút gọn \verb|git status --short| rất tiện khi làm bài: hai cột ký tự đầu dòng thể hiện trạng thái của index và của thư mục làm việc.
Ví dụ \verb|??| nghĩa là untracked, \verb|A| nghĩa là đã add, \verb|M| nghĩa là đã sửa đổi.
Quy trình chuẩn cho commit đầu tiên của project ASPX gồm bốn bước trong Khối lệnh dưới đây.

\begin{lstlisting}[language=bash]
# 1. Init a Git repo inside the ASPX project folder
git init AspxDemo_NET481
cd AspxDemo_NET481

# 2. Local identity for this repo only
git config user.name "Nguyen Van A"
git config user.email "nva@student.edu.vn"

# 3. Create a simple ASPX page, then check status
echo '<%@ Page Language="C#" %>' > Default.aspx
echo '<html><body>Hello ASPX</body></html>' >> Default.aspx
git status --short
# Output: ?? Default.aspx   (untracked)

# 4. Stage the file, then commit
git add Default.aspx
git status --short
# Output: A  Default.aspx   (staged)

git commit -m "feat: add Default.aspx homepage skeleton"
git log --oneline -1

# 5. Edit the file -> modified (unstaged)
echo '<!-- header -->' >> Default.aspx
git status --short
# Output:  M Default.aspx   (modified)

# 6. Forgot a file? Stage it and fold into the last commit
echo '<%@ Master Language="C#" %>' > Site.Master
git add Default.aspx Site.Master
git commit --amend --no-edit
git log --oneline -1
git status --short
\end{lstlisting}

Tùy chọn \verb|--amend| viết lại commit mới nhất bằng cách thay cây file và thông điệp của nó, thay vì tạo commit mới.
Lệnh này chỉ an toàn khi commit chưa được push lên kho chứa chung, vì nó thay đổi mã băm của commit \cite{progit}.
Với project bài tập, nhóm nên viết thông điệp commit theo chuẩn Conventional Commits, ví dụ \verb|feat: ...|, \verb|fix: ...|, \verb|docs: ...| \cite{convcommits}.
Thông điệp tốt giúp giảng viên đọc lịch sử là hiểu ngay tiến độ: trang nào đã xong, lỗi nào đã sửa, tài liệu nào đã cập nhật.
Một quy ước gợi ý cho đồ án Web Forms là mỗi commit chỉ đụng một trang hoặc một UserControl, kèm thông điệp nêu rõ tên file đã đổi.

\subsection{Xem lịch sử và so sánh}
\label{subsec:c22-history}

Ba nhu cầu phổ biến nhất khi đọc lịch sử là xem danh sách commit, xem chi tiết một commit và so sánh các phiên bản với nhau.
Lệnh \verb|git log| hiển thị lịch sử, lệnh \verb|git show| hiển thị chi tiết, còn lệnh \verb|git diff| hiển thị sự khác biệt giữa hai trạng thái \cite{progit}.
Bảng~\ref{tab:c22-log} tóm tắt các tùy chọn mà sinh viên dùng hằng ngày khi làm bài tập nhóm.

\begin{table}[H]
\centering
\caption{Các lệnh xem lịch sử và so sánh dùng hằng ngày}
\label{tab:c22-log}
\begin{tabular}{lll}
\toprule
Mục đích & Lệnh & Ghi chú \\
\midrule
Lịch sử rút gọn & \texttt{git log --oneline} & Mỗi commit một dòng \\
Sơ đồ nhánh & \texttt{git log --graph --all} & Nên kèm \texttt{--oneline --decorate} \\
Chi tiết một commit & \texttt{git show} + mã băm & Gồm metadata và patch \\
So sánh chưa stage & \texttt{git diff} & Thư mục làm việc so với index \\
So sánh đã stage & \texttt{git diff --staged} & Index so với HEAD \\
Thống kê file đổi & \texttt{git diff --stat} & Gọn, dùng khi review \\
\bottomrule
\end{tabular}
\end{table}

Tùy chọn \verb|--oneline| nén mỗi commit thành một dòng gồm bảy ký tự đầu của mã băm và dòng tiêu đề của thông điệp.
Kết hợp ba tùy chọn \verb|log --oneline --graph --all| cho ra sơ đồ ASCII thể hiện đầy đủ các nhánh và điểm rẽ, điểm hợp nhất.
Lệnh \verb|git show| không tham số sẽ hiển thị commit mới nhất, còn khi kèm mã băm thì hiển thị đúng commit đó.
Thêm tên file ở cuối, ví dụ \verb|git show a1b2c3d -- Default.aspx|, để chỉ xem thay đổi trên file quan tâm.

\begin{lstlisting}[language=bash]
# Compact history, newest first
git log --oneline -5
# a1b2c3d feat: add login form validation
# 9f8e7d6 fix: correct Web.config debug flag
# 4c3b2a1 docs: update README build steps
# 77aa11b feat: add cart UserControl
# 2e1f0c9 init: scaffold ASPX Web Forms project

# Graph of all branches with decorations
git log --oneline --graph --all --decorate -8
# * a1b2c3d (HEAD -> main) feat: add login form validation
# | * 88bb22c (feature/cart) feat: add cart UserControl
# |/
# * 9f8e7d6 fix: correct Web.config debug flag
# * 4c3b2a1 docs: update README build steps

# Show one commit in detail (metadata + patch)
git show --stat HEAD
git show a1b2c3d -- Default.aspx
\end{lstlisting}

Lệnh \verb|git diff| không tham số so sánh thư mục làm việc với vùng stage, tức chỉ hiện thay đổi chưa được add.
Muốn xem thay đổi đã add nhưng chưa commit, bắt buộc phải thêm \verb|--staged| (hoặc \verb|--cached|, hai tên tương đương).
Ký hiệu hai chấm \verb|main..feature/cart| so sánh đầu nhánh này với đầu nhánh kia, rất tiện trước khi hợp nhất hoặc nộp pull request.
Đoạn patch trong output dùng dấu cộng cho dòng thêm và dấu trừ cho dòng xóa, kèm tiêu đề hunk dạng \verb|@@ -1,2 +1,3 @@|.

\begin{lstlisting}[language=bash]
# Unstaged changes: working tree vs index
git diff
# diff --git a/Default.aspx b/Default.aspx
# --- a/Default.aspx
# +++ b/Default.aspx
# @@ -1,2 +1,3 @@
#  <%@ Page Language="C#" %>
# +<!-- header -->
#  <html><body>Hello ASPX</body></html>

# Staged changes: index vs HEAD
git add Default.aspx
git diff --staged
git diff --cached --stat

# Compare branches or commits before merge or review
git diff main..feature/cart --stat
git diff HEAD~1 HEAD -- Web.config
git diff --word-diff Default.aspx
\end{lstlisting}

Một thói quen tốt là luôn chạy \verb|git diff --staged| để rà soát lần cuối trước khi commit, giống như đọc lại bài trước khi nộp.
Với file \verb|Web.config| chứa chuỗi kết nối, bước rà soát này giúp phát hiện sớm mật khẩu thật bị lọt vào commit.
Sinh viên nên tập đọc patch thành thạo vì giao diện review trên GitHub và GitLab cũng hiển thị đúng định dạng này \cite{githubdocs}.
\subsection{Hoàn tác thay đổi}
\label{subsec:c22-undo}

Git phân biệt ba tầng có thể hoàn tác: commit đã tạo, vùng stage và thư mục làm việc.
Năm lệnh hoàn tác phổ biến gồm ba chế độ của \verb|git reset|, cộng thêm \verb|git revert| và \verb|git restore|.
Bảng~\ref{tab:c22-undo} so sánh tác động của từng lệnh để sinh viên chọn đúng công cụ cho từng tình huống \cite{progit}.

\begin{table}[H]
\centering
\caption{So sánh các lệnh hoàn tác thay đổi trong Git}
\label{tab:c22-undo}
\begin{tabular}{llll}
\toprule
Lệnh & Tác động & An toàn trên nhánh chung & Khi dùng \\
\midrule
\texttt{reset --soft} & Dời HEAD, giữ index và file & Không & Gộp commit cục bộ chưa push \\
\texttt{reset --mixed} & Dời HEAD và index, giữ file & Không & Bỏ stage hàng loạt \\
\texttt{reset --hard} & Dời cả ba, mất sửa đổi & Không, xem cảnh báo & Xóa sạch nháp cục bộ \\
\texttt{revert} & Tạo commit mới đảo ngược & Có & Hoàn tác commit đã push \\
\texttt{restore} & Chỉ sửa file, không đụng lịch sử & Có & Hủy sửa chưa commit \\
\bottomrule
\end{tabular}
\end{table}

\textbf{Cảnh báo:} tuyệt đối không dùng \verb|git reset --hard| hay \verb|git push --force| trên nhánh chung như \verb|main| hoặc \verb|develop|.
Hai lệnh này viết lại lịch sử công khai, khiến đồng đội bị mất commit hoặc phải xử lý xung đột không đáng có \cite{atlassiangit}.
Quy tắc ngón tay cái là lịch sử đã push thì chỉ được sửa bằng commit mới, tức dùng \verb|git revert| thay vì \verb|reset|.
Ngoại lệ duy nhất cho \verb|--hard| là dọn dẹp nháp hoàn toàn cục bộ, khi chắc chắn không ai khác có các commit đó.
Trước mọi lệnh nguy hiểm, hãy tạo nhánh dự phòng bằng \verb|git branch backup-ten-ngay| để luôn có đường quay lại.

\begin{lstlisting}[language=bash]
# --- git restore: undo files, history untouched ---
git restore Default.aspx            # discard unstaged edits
git restore --staged Default.aspx   # unstage, keep edits in files
git restore --source=HEAD~1 Web.config  # take file from old commit

# --- git reset: move HEAD (local drafts only) ---
git reset --soft HEAD~1    # keep index + working tree
git commit -m "feat: squash login + validation into one commit"
git reset --mixed HEAD~1   # keep working tree, reset index
git reset HEAD Default.aspx  # unstage one file (mixed is default)

# --- git revert: safe undo on shared branches ---
git log --oneline -3
git revert HEAD --no-edit        # new commit that undoes HEAD
git revert --no-commit 9f8e7d6  # stage a revert, commit later
git status --short

# --- last resort: discard everything local (see warning) ---
git branch backup-2026-05-30     # safety branch first
git reset --hard HEAD            # drop all uncommitted changes
\end{lstlisting}

Lệnh \verb|git restore| ra đời để tách bạch hai việc mà trước đây đều dồn vào \verb|git checkout|: lấy lại file và chuyển nhánh.
Cú pháp \verb|git restore --source| cho phép lấy nội dung file từ bất kỳ commit nào mà không di chuyển HEAD.
Lệnh \verb|git revert| không xóa commit cũ mà tạo commit mới có nội dung đảo ngược, nên lịch sử vẫn trung thực và đầy đủ.
Khi revert một commit hợp nhất, cần thêm \verb|-m 1| để chỉ rõ giữ lại phía nhánh chính \cite{progit}.
Với sinh viên mới học, thứ tự ưu tiên khi lỡ tay là \verb|restore| trước, \verb|revert| sau, và chỉ dùng \verb|reset| khi hiểu rõ tác động.
\subsection{Stash và tag}
\label{subsec:c22-stash-tag}

Khi đang sửa dở một trang mà phải chuyển nhánh gấp để sửa lỗi nóng, lệnh \verb|git stash| cất tạm mọi thay đổi chưa commit vào ngăn xếp.
Vùng làm việc trở lại sạch sẽ ngay lập tức, và sinh viên có thể lấy lại công việc dở dang bất cứ lúc nào bằng \verb|git stash pop| \cite{progit}.
Ba thao tác dùng nhiều nhất là \verb|stash push| để cất, \verb|stash list| để liệt kê và \verb|stash pop| để lấy lại đồng thời xóa khỏi ngăn xếp.
Nên kèm thông điệp bằng tùy chọn \verb|-m|, nếu không Git sẽ đặt tên mặc định khó phân biệt khi có nhiều mục stash.

\begin{lstlisting}[language=bash]
# Stash unfinished work before switching branch
git status --short
#  M Default.aspx
#  M Login.aspx
git stash push -m "wip: half-done login check"
git status --short   # clean again
git stash list
# stash@{0}: wip: half-done login check

# Fix the hot bug on another branch, then come back
git switch hotfix-config
# ... edit Web.config, commit ...
git switch main
git stash pop        # restore + drop stash@{0}
git status --short
git stash list
\end{lstlisting}

Cần phân biệt \verb|stash pop| với \verb|stash apply|: lệnh \verb|pop| xóa mục stash sau khi lấy lại, còn \verb|apply| giữ lại bản copy.
Nếu ngăn xếp có nhiều mục, chỉ rõ mục cần lấy bằng tên dạng \verb|stash@{0}| hoặc \verb|stash@{1}| ở cuối lệnh.
Trường hợp stash chứa file mới hoàn toàn (untracked), phải thêm tùy chọn \verb|-u| thì Git mới cất cả các file đó.

Trong khi stash quản lý công việc tạm thời, tag dùng để đánh dấu các mốc quan trọng như bản nộp bài hoặc bản phát hành.
Git có hai loại tag: tag nhẹ (lightweight) chỉ là một tên trỏ tới commit, còn tag annotated là một đối tượng đầy đủ có tác giả, ngày tháng và thông điệp \cite{progit}.
Với bài nộp môn học, nhóm nên dùng tag annotated để giảng viên kiểm tra được ai đánh tag và nộp những gì.
Quy ước tên gợi ý là \verb|v1.0|, \verb|v1.1| cho các lần nộp, kèm tiền tố lớp nếu giảng viên yêu cầu.

\begin{lstlisting}[language=bash]
# Lightweight tag: just a name on a commit
git tag v1.0-btvn

# Annotated tag: recommended for course submissions
git tag -a v1.0 -m "Nop bai lan 1: Web Forms chay duoc"
git show v1.0 --no-patch   # author, date, message
git tag -l "v1.*"          # list matching tags
git tag -n99               # list tags with messages

# Share tags with the remote (tags are NOT pushed by default)
git push origin v1.0       # push one tag
git push origin --tags     # push all tags
git ls-remote --tags origin
\end{lstlisting}

Lưu ý quan trọng là lệnh \verb|git push| mặc định không đẩy tag lên kho chứa chung, nên phải push tag một cách tường minh.
Muốn xóa tag nhầm tên, dùng \verb|git tag -d v1.0| cho tag cục bộ rồi \verb|git push origin :v1.0| để xóa tag trên remote.
Trước giờ nộp bài, cả nhóm nên thống nhất một commit chốt trên nhánh \verb|main|, đánh tag annotated và kiểm tra lại bằng \verb|git show|.
Việc kết hợp stash cho công việc dở dang và tag cho các mốc nộp bài giúp quy trình nhóm ba người luôn gọn gàng và dễ truy vết.
